\section{Introducción}

El objetivo principal de este proyecto es diseñar e implementar una \textbf{base de datos} musical
para archivos de audio en formato \textbf{MP3} con etiquetas \textbf{ID3v2.4}. El sistema cuenta con
al menos 4 módulos bien definidos: un \textbf{minero}, una \textbf{base de datos}, un \textbf{lenguaje de consultas}
y una \textbf{interfaz gráfica}.

\section*{Alcance}

\textbf{8-Track} contempla cubrir el flujo de la especificación del \textbf{Proyecto 2}, es decir:
escaneo de directorios con archivos MP3, base de datos con la biblioteca musical, edición de
metadatos en la base de datos, búsquedas avanzadas a partir de un \textbf{DSL} y una intefaz
gráfica que permita realizar todas las \textit{features} de la especificación.

Adicionalmente, se contempla extender las funcionalidades básicas mediante la incorporación de
reproducción de audio directamente desde la aplicación, así como el diseño de una interfaz intuitiva,
visualmente atractiva y orientada a ofrecer una experiencia de usuario fluida y enriquecida.

\section*{Limitaciones}

La principal limitación del proyecto es su compatibilidad restringida con formatos de audio, pues
únicamente admite archivos \textbf{MP3}. En consecuencia, las colecciones musicales heterogéneas que
incluyan formatos populares como \textbf{FLAC}, \textbf{ALAC}, \textbf{WAV} u \textbf{OGG} no podrán administrarse en su totalidad.
Por ello, la incorporación de soporte para formatos adicionales representa una de las principales
oportunidades de expansión del sistema.

Además, la aplicación opera bajo el principio de \textit{Garbage In, Garbage Out} (GIGO), pues la calidad y
precisión de la información presentada dependen directamente de los \textbf{metadatos} contenidos
en los archivos \textbf{MP3}. En consecuencia, elementos como los nombres de álbumes, las portadas y
la información de los artistas pueden resultar incompletos, incorrectos o incluso inexistentes si
los archivos originales carecen de los datos correspondientes.

Esta dependencia podría mitigarse mediante la integración de \textbf{APIs de terceros}, como las proporcionadas
por \textbf{MusicBrainz}, \textbf{Last.fm} o \textbf{Deezer}, que permitirían consultar fuentes externas
para complementar, verificar y enriquecer los metadatos de la biblioteca musical, así como recuperar
portadas de álbumes e imágenes de artistas. Sin embargo, estas integraciones quedan fuera del alcance
actual del proyecto y se consideran posibles extensiones futuras.

El proyecto está completamente orientado a usuarios de \textbf{Linux} y no contempla una extensión
multiplataforma.
